% Auto-generated by scripts/make_hypothesis_rationale_table.R --- do not edit by hand.
% Body rows for the main-body longtable defined in draft_APSG.tex
% (label = tab:hypotheses-rationale).
01 & topic-avoidance & Countries with lower Liberal democracy levels will exhibit lower mean PCI scores in SSH abstracts. & Researchers working under autocratic rule face institutional incentives to avoid topics that may attract state scrutiny, including explicitly political subject matter. \\[3pt]
02 & topic-avoidance & Countries with lower Liberal democracy levels will exhibit lower Shannon entropy of author-supplied keyword distributions at the country-year level. & Author-supplied keywords are the researcher's own signal of what a paper is about; in autocratic settings, researchers who self-censor will gravitate toward safer, institutionally approved topics, causing keyword vocabularies to concentrate around a narrower range of terms. \\[3pt]
03 & topic-avoidance & Countries with lower Liberal democracy levels will produce a lower share of SSH output in politically sensitive disciplines. & Self-censorship does not operate only within disciplines — it also operates across them. \\[3pt]
04 & topic-avoidance & Countries with lower Liberal democracy levels will have a lower share of SSH articles containing regime-sensitive terms in titles, author keywords, and abstracts. & Self-censorship theory predicts that researchers in autocracies avoid topics that could draw state scrutiny or threaten their careers. \\[3pt]
05 & topic-avoidance & Countries with lower Liberal democracy levels will exhibit a lower share of SSH articles critically examining their own governance and institutions. & Self-censorship theories predict that researchers operating under authoritarian constraints avoid producing work that could be perceived as threatening to the regime. \\[3pt]
06 & topic-avoidance & Countries with lower Liberal democracy levels will exhibit lower mean pairwise semantic diversity of SSH abstracts within country-field-year clusters. & If researchers in autocracies self-censor by converging on politically safe topics and framings, the full semantic space of their published output — not just the presence or absence of specific keywords — should narrow. \\[3pt]
07 & topic-avoidance & Countries with lower Liberal democracy levels will have a lower share of SSH articles acknowledging international funding agencies. & Receiving funding from international agencies — such as the US National Science Foundation, the European Research Council, the World Bank, the Ford Foundation, or bilateral aid bodies — involves agreeing to donor norms, submitting to external review, and producing work that is visible to international audiences. \\[3pt]
08 & framing-neutrality & Countries with lower Liberal democracy levels will exhibit higher mean epistemic hedging rates in SSH abstracts. & Researchers working under autocratic constraints face institutional incentives not only to avoid politically sensitive topics outright, but also to hedge and soften any claims they do make, deflecting regime scrutiny by rendering findings provisional rather than declarative. \\[3pt]
09 & framing-neutrality & Countries with lower Liberal democracy levels will exhibit lower mean normative language rates in SSH abstracts. & Self-censorship theory predicts that scholars in autocracies strategically minimise language that could be read as a political or moral claim against the regime. \\[3pt]
10 & framing-neutrality & Countries with lower Liberal democracy levels will exhibit a higher share of SSH abstracts classified as purely technocratic (no political references). & Self-censorship theory predicts that researchers in autocracies will actively reframe their work to minimize the risk of regime scrutiny — not only by avoiding sensitive topics altogether, but by stripping politically referential language from work they do publish, presenting it instead in purely technical or administrative terms. \\[3pt]
11 & framing-neutrality & Countries with lower Liberal democracy levels will exhibit lower mean rates of first-person argumentative stance phrases in SSH abstracts. & Making an explicit attributed claim in academic writing ("we argue that X") exposes the author as the source of an assertion, increasing personal accountability for the content of a paper. \\[3pt]
12 & framing-neutrality & Countries with lower Liberal democracy levels will exhibit a lower share of SSH abstracts ending with an explicit normative or policy conclusion. & Self-censorship theory predicts that researchers under autocratic constraints avoid signaling political engagement. \\[3pt]
13 & framing-neutrality & Countries with lower Liberal democracy levels will publish a higher share of SSH output in domestically dominated journals. & Publishing in an international journal subjects a paper to editorial standards, reviewer networks, and professional norms anchored in liberal academic culture; publishing in a domestically dominated journal places the work in a locally controlled institutional environment where regime-aligned editorial gatekeepers have greater influence. \\[3pt]
14 & collaboration-constraint & Countries with lower Liberal democracy levels will exhibit a lower mean number of distinct co-author countries per SSH article. & International collaboration in research is not politically neutral: maintaining professional contacts abroad, exchanging manuscripts across borders, and attending international conferences all expose researchers to cross-border oversight and potential regime scrutiny. \\[3pt]
15 & collaboration-constraint & Countries with lower Liberal democracy levels will have a lower share of SSH articles co-authored with researchers from liberal democracies (v2x\_libdem > 0.5). & Collaboration with scholars in liberal democracies is not politically neutral for researchers in autocracies: democratic co-authors operate under editorial and professional norms that may encourage politically sensitive research agendas, and democratic-country affiliations may trigger regime suspicion toward the collaboration and its outputs. \\[3pt]
16 & collaboration-constraint & Countries with lower Liberal democracy levels will exhibit a higher share of SSH articles with exclusively domestic authorship. & Self-censorship theory predicts that researchers in autocracies constrain their collaboration choices to minimize exposure to foreign oversight, documentation of foreign contacts, and the transmission of liberal academic norms. \\[3pt]
17 & collaboration-constraint & Countries with lower Liberal democracy levels will exhibit a lower mean number of authors per SSH article. & Collaborative research requires multiple people to share ideas, access data, and jointly produce work that may attract political scrutiny. \\[3pt]
18 & collaboration-constraint & Countries with lower Liberal democracy levels will exhibit a lower mean number of distinct author institutions per SSH article. & Inter-institutional collaboration requires spreading a research project across organizational boundaries — across universities, research centers, government agencies, and other bodies. \\[3pt]
19 & visibility-suppression & Countries with lower Liberal democracy levels will exhibit a lower mean normalized citation ratio for SSH articles. & Self-censorship produces work that is rhetorically cautious, topically narrow, and unlikely to engage the politically salient debates that drive citation traffic in SSH; hedged claims and depoliticized framings give other scholars little to build on, criticize, or extend, so on average autocracy-produced SSH work should attract fewer citations than comparable work from democratic contexts. \\[3pt]
20 & visibility-suppression & The citation penalty associated with lower Liberal democracy levels will be larger for politically sensitive articles than for non-sensitive articles, as indicated by a positive coefficient on the v2x\_libdem × sensitive\_flag interaction term. & If self-censorship operates through cautious framing and weakened claims, the citation penalty should not apply uniformly: politically non-sensitive SSH topics impose few framing constraints in either regime context, so citations should be comparable, while sensitive topics force autocratic authors to hedge, depoliticize, or narrow their contributions in ways their democratic counterparts are not constrained to do --- producing a citation gap that concentrates precisely where the theory predicts self-censorship bites hardest. \\[3pt]
21 & visibility-suppression & Countries with lower Liberal democracy levels will exhibit a lower share of SSH articles that receive at least one citation within a country-year. & Articles that fail to advance any distinctive claim are the ones most likely to attract zero citations; if self-censorship pushes researchers toward incremental, low-stakes contributions on safe topics, autocratic SSH output should exhibit a larger share of articles that never enter scholarly conversation at all. This is the extensive-margin counterpart to the average-impact test in Team 19. \\[3pt]
22 & visibility-suppression & Countries with lower Liberal democracy levels will exhibit a higher Gini coefficient of citations across SSH articles. & Self-censorship may bifurcate autocratic SSH output: a small set of articles --- those touching on regime-endorsed priorities, applied development topics, or internationally co-funded projects --- escape the constraints and accumulate citations normally, while the conformist majority is undifferentiated and rarely cited. Citation distribution within autocratic country-years should therefore be more skewed than in democratic ones, captured by a higher within-cell citation Gini. \\[3pt]
23 & ideological-alignment & Countries with lower Liberal democracy levels will exhibit a lower share of SSH abstracts with liberal/neutral political-economy framing (NEUTRAL + NONE) and a higher share with ideologically aligned framing (LEFT or RIGHT-CONSERVATIVE). & Self-censorship under autocracy is not ideologically neutral: researchers in communist regimes face pressure to produce scholarship consonant with statist, collectivist, and anti-market frames, while researchers in nationalist or right-authoritarian regimes face complementary pressures toward market-nationalist or traditionalist frames. \\[3pt]
24 & ideological-alignment & Countries with lower Liberal democracy levels will exhibit a higher share of SSH abstracts that actively legitimize or endorse the current political system, state authority, or leadership. & Self-censorship theories have largely focused on what researchers avoid; yet autocratic regimes do not merely suppress dissent — they actively incentivize, fund, and reward scholarship that validates state power and ideology. \\[3pt]
25 & ideological-alignment & Countries with lower Liberal democracy levels will exhibit a higher share of SSH abstracts that actively frame liberal democracy, Western political institutions, or international human rights norms as fundamentally flawed or illegitimate. & Self-censorship theory predicts that researchers in autocracies avoid politically sensitive topics to minimize career risk, but the same logic of regime alignment generates an additional prediction: where states actively promote counter-liberal ideologies, researchers face positive incentives to produce scholarship consonant with those ideologies. \\[3pt]
